\documentclass[12pt,a4paper]{article}
\usepackage[utf8]{inputenc}
\usepackage[T1]{fontenc}
\usepackage{amsmath}
\usepackage{booktabs}
\usepackage{geometry}
\usepackage{hyperref}
\usepackage{cite}
\geometry{a4paper,top=2.5cm,bottom=2.5cm,left=2.5cm,right=2.5cm}
\hypersetup{colorlinks=true,linkcolor=blue,citecolor=blue,urlcolor=blue}

\title{\textbf{Compact Spherical Reactor Design for\\
Proton--Electron Capture Chain Energy Extraction:\\
Multi-Layer Architecture with Thermoelectric Output}}
\author{Eymen Aydoğan}
\date{May 19, 2026}

\begin{document}
\maketitle

\begin{abstract}
We present a compact spherical reactor design implementing the PECCNAL
proton--electron capture chain reaction for direct electricity
generation. The design employs a multi-layer architecture: a central
reaction chamber containing closely-spaced nickel electrode arrays
in deuterium oxide (D$_2$O) with dissolved Lithium-6, surrounded by
a copper thermal conductor layer, and an outer array of thermoelectric
(Peltier) modules converting thermal output to electricity.
The electrode array consists of alternating nickel and stainless
steel plates spaced 1--2 mm apart, maximizing reaction surface area
while minimizing proton travel distance. Preliminary scaling analysis
suggests a 15~cm diameter prototype could produce measurable
electrical output detectable with standard instrumentation.
\end{abstract}

\noindent\textbf{Keywords:} compact reactor, thermoelectric,
peltier, nickel electrodes, spherical design, D$_2$O, prototype

\hrule\vspace{1em}

\section{Introduction}
Previous papers in this series have established the theoretical
basis for PECCNAL reactions. Here we present a concrete engineering
design for a compact prototype reactor intended for laboratory
demonstration and calorimetric verification.

\section{Reactor Architecture}

\subsection{Layer Structure}

\begin{table}[h!]
\centering
\caption{Reactor layer specifications}
\vspace{0.5em}
\begin{tabular}{llll}
\toprule
\textbf{Layer} & \textbf{Material} & \textbf{Function} & \textbf{Thickness} \\
\midrule
Core & Ni/SS electrodes + D$_2$O + Li$_2$CO$_3$ & Reaction & 60 mm dia. \\
Thermal & Copper & Heat conduction & 5 mm \\
Insulation & PETG polymer & Electrical isolation & 2 mm \\
Output & TEC1-12706 Peltier array & Heat→electricity & 4 mm \\
Shield & Polyethylene & Neutron absorption & 20 mm \\
Housing & Aluminium & Structural & 3 mm \\
\bottomrule
\end{tabular}
\end{table}

Total diameter: $\approx$ 190 mm. Total mass: $\approx$ 2.5 kg.

\subsection{Electrode Array}

The electrode array consists of:
\begin{itemize}
\item 50 nickel plates (cathode, $-$), 60 mm $\times$ 60 mm $\times$ 0.5 mm
\item 50 stainless steel plates (anode, $+$), same dimensions
\item Plate spacing: 1 mm
\item Total electrode surface area: $50 \times 2 \times 36\,\text{cm}^2 = 3600\,\text{cm}^2$
\item PETG polymer spacers (3D-printed) maintaining alignment
\end{itemize}

\subsection{Electrolyte}

D$_2$O with 5\% dissolved Li$_2$CO$_3$ by mass:
\begin{itemize}
\item Volume: $\approx$ 100 mL
\item Li-6 content: $7.5\%$ of Li mass (natural abundance)
\item Conductivity enhanced by Li$_2$CO$_3$ dissociation
\end{itemize}

\section{Energy Flow}

\begin{equation}
\text{Electrical input} \xrightarrow{\text{electrolysis}}
\text{D}^+ \text{ ions} \xrightarrow{\text{PECCNAL}}
\text{Heat} \xrightarrow{\text{Peltier}}
\text{Electrical output}
\end{equation}

Peltier efficiency at $\Delta T = 10$ K:
\begin{equation}
\eta_\text{Peltier} \approx 5\text{--}8\%
\end{equation}

\section{Measurement Protocol}

\begin{enumerate}
\item Measure input power: $P_\text{in} = V \times I$ (multimeter)
\item Measure output power: $P_\text{out} = V_\text{Peltier} \times I_\text{Peltier}$
\item Control experiment: replace D$_2$O with H$_2$O
\item Excess power: $\Delta P = P_\text{out,D_2O} - P_\text{out,H_2O}$
\item If $\Delta P > 0$ consistently: PECCNAL reaction confirmed
\end{enumerate}

\section{Manufacturing}

\begin{itemize}
\item Electrode spacers: 3D-printed PETG (Elegoo Neptune 4)
\item Copper housing: cut from sheet stock, soldered
\item Peltier modules: TEC1-12706, commercially available
\item Total estimated cost: 200--400 USD
\end{itemize}

\section{Conclusion}
A compact 190 mm spherical reactor implementing the PECCNAL
mechanism is described, using 100 alternating Ni/SS electrode
plates in D$_2$O/Li$_2$CO$_3$ electrolyte. The design is
manufacturable with 3D printing and standard electronic
components at low cost. Thermoelectric output and calorimetric
excess heat measurement provide two independent verification
channels for the PECCNAL reaction.

\begin{thebibliography}{9}
\bibitem{aydogan2026a} E. Aydoğan, ``Coulomb-Barrier-Free Neutron Production,'' arXiv preprint, 2026.
\bibitem{aydogan2026b} E. Aydoğan, ``Deuterium-Mediated Neutron Chain Fusion,'' arXiv preprint, 2026.
\bibitem{aydogan2026c} E. Aydoğan, ``Self-Sustaining Tritium Regeneration Loop,'' arXiv preprint, 2026.
\bibitem{aydogan2026e} E. Aydoğan, ``Heavy Water Electrolysis with Nickel Electrodes,'' arXiv preprint, 2026.
\bibitem{pdg2022} R. L. Workman et al. (PDG), \textit{Prog. Theor. Exp. Phys.}, 083C01, 2022.
\end{thebibliography}
\end{document}
